% Generated by `taskgen course examples` from number_theory.inverse (seed 7).
% Do not edit: change examples.yaml and regenerate.
\begin{example}[обратный элемент по составному модулю]\label{ex:euclid-inverse}
Найти элемент, обратный к $a = 33$ в кольце вычетов по модулю $n = 146$, или доказать, что его не существует.

\par\noindent\textit{Решение.} Обратный к $a$ существует тогда и только тогда, когда $\gcd(a, n) = 1$; расширенный алгоритм Евклида для $(n, a)$ даёт $d = \gcd(n, a)$ и $u n + v a = d$.\[\begin{array}{r|r|r|r}
r_i & q_i & u_i & v_i \\ \hline
146 &  & 1 & 0 \\
33 &  & 0 & 1 \\
14 & 4 & 1 & -4 \\
5 & 2 & -2 & 9 \\
4 & 2 & 5 & -22 \\
1 & 1 & -7 & 31 \\
0 & 4 & 33 & -146 \\
\end{array}\qquad r_i = u_i \cdot 146 + v_i \cdot 33\]
$d = 1 = -7 \cdot 146 + 31 \cdot 33$, поэтому $33^{-1} \equiv 31\ (\textrm{mod}\ 146)$ (столбец $u_i$ можно было не вычислять).\par Проверка: $33 \cdot 31 = 1023 = 7 \cdot 146 + 1$.\[33^{-1} \equiv 31\ (\textrm{mod}\ 146).\]
\end{example}
